\ResumenAbstract
{%
La presente investigación desarrolló un sistema web automatizado para optimizar la rotación de circuitos eléctricos ante situaciones de déficit de generación en la red de distribución del sistema electroenergético cubano. Su objetivo principal fue sustituir el proceso manual actual por una herramienta computacional de arquitectura cliente-servidor que analizara múltiples variables en tiempo real —como cargas horarias, MWh totales por circuito, averías, aseguramientos y trabajos planificados— para generar órdenes de afectación y restablecimiento óptimas siguiendo directrices del Despacho Nacional de Carga. La metodología combinó el análisis de requerimientos con el diseño e implementación de una aplicación web moderna utilizando el framework Next.js y la biblioteca React para el frontend, estilizado con Tailwind CSS, y un backend robusto en NestJS que integró bases de datos SQL Server 2008. La gestión de estados se centralizó mediante Hooks y Context API, permitiendo una interfaz dinámica y eficiente. Los principales hallazgos evidenciaron una reducción significativa en el tiempo de toma de decisiones (de 30-45 minutos a 2-3 segundos), mayor precisión en la selección de circuitos y la capacidad de simular escenarios de demanda. Se concluyó que la herramienta incrementó la eficiencia operativa del Despacho Provincial de Carga y validó la transición de sistemas de escritorio a plataformas web escalables bajo metodologías ágiles. \noindent\textbf{PALABRAS CLAVE}\newline Automatización eléctrica, Rotación de circuitos, Next.js, NestJS, Programación Extrema (XP).
}
{%
This research developed an automated web-based system to optimize the rotation of electrical circuits during generation deficit situations in the distribution network of the Cuban power system. Its main objective was to replace the current manual process with a client-server computational tool that analyzed multiple variables in real time---such as hourly loads, total MWh per circuit, failures, secured circuits, and planned maintenance---to generate optimal shutdown and restoration orders following guidelines from the National Load Dispatch Center. The methodology combined requirements analysis with the design and implementation of a modern web application using the Next.js framework and React library for the frontend, styled with Tailwind CSS, and a robust NestJS backend integrating SQL Server 2008 databases. State management was centralized through Hooks and Context API, enabling a dynamic and efficient interface. Main findings showed a significant reduction in decision-making time (from 30-45 minutes to 2-3 seconds), greater accuracy in circuit selection, and the ability to simulate demand scenarios. It was concluded that the tool increased the operational efficiency of the Provincial Load Dispatch Center and validated the transition from desktop systems to scalable web platforms under agile methodologies. \noindent\textbf{KEYWORDS}\newline Electrical automation, Circuit rotation, Next.js, NestJS, Extreme Programming (XP).
}